\subsection{An explicit nonzero commutator in finite geometric quotients}
\label{subsec:cm-low-class-two}

This computation specifies a small nonabelian quotient of the
outer image.  Its definition does not assert detection of the full
class-two quotient of the arithmetic source.
Let \(\Pi=\pi_1^{(p)}((E-\{O\})_{\bar K})\), with descending central
series \(\Pi(n)\), and put
\[
 P=\Pi/\Pi(13),\qquad
 P_r=P/\overline{\langle x^{p^r}:x\in P\rangle}\quad(r\geq1).
\]
These are characteristic quotients of \(\Pi\). The successive
lower-central quotients of \(P_r\) are finitely generated abelian
groups of exponent dividing \(p^r\), so every \(P_r\) is a finite \(p\)-group.
The field identity~\eqref{eq:cm-radical-outer-field} gives the
actual induced outer map \(\rho_H:H\longrightarrow\operatorname{Out}(\Pi)\):
indeed, restriction identifies
\(\operatorname{Gal}(K(E[p])\Omega/K(E[p]))\) with
\(\operatorname{Gal}(\Omega/L_1)\), because the intersection is
\(L_1\), and the outer action factors through its quotient
\(\operatorname{Gal}(\Omega^*/L_1)\).
Define the finite image
\[
 Q_r=\operatorname{im}\bigl(H\xrightarrow{\rho_H}
                              \operatorname{Out}(P_r)\bigr).
\]

\begin{proposition}\label{prop:cm-small-nonabelian-image}
For every class-number-one imaginary quadratic \(K\), every
split \(p\geq5\), and every \(E/K\) as above, all \(Q_r\) have
nilpotency class at most two.  They are nonabelian for all
sufficiently large \(r\).
\end{proposition}

\begin{proof}
We first give the rational Lie calculation.
Let \(L\) be the free Lie algebra over \(\mathbb Q\) on \(a,b\).
Let \(D_{rs}\) denote its space of derivations with
\[
 D(a)\in L_{r+1,s},\qquad D(b)\in L_{r,s+1},\qquad
                         D([a,b])=0.
\]
For a Lyndon word \(w\), write \([w]\) for its standard bracket:
split off its longest proper Lyndon suffix and bracket
recursively, with \(a<b\).
The following coefficient vectors define \(d\in D_{22}\) and
\(e,f\in D_{33}\).  The entries of each row are coefficients in
the displayed ordered Lyndon basis:
\[
\begin{array}{c|c|r@{\,,\,}r@{\,,\,}r@{\,,\,}r@{\,,\,}r}
 &\text{ordered basis}&\multicolumn{5}{c}{\text{coefficients}}\\ \hline
d(a)&[aaabb],[aabab]&\multicolumn{5}{c}{(-2,-1)}\\
d(b)&[aabbb],[ababb]&\multicolumn{5}{c}{(2,3)}\\
e(a)&[aaaabbb],[aaababb],[aaabbab],[aabaabb],[aababab]
       &1&-5&-6&1&-2\\
e(b)&[aaabbbb],[aababbb],[aabbabb],[aabbbab],[abababb]
       &-1&4&6&8&0\\
f(a)&[aaaabbb],[aaababb],[aaabbab],[aabaabb],[aababab]
       &-3&-9&-6&-3&-2\\
f(b)&[aaabbbb],[aababbb],[aabbabb],[aabbbab],[abababb]
       &3&12&6&0&8
\end{array}
\]
Direct expansion gives
\(d([a,b])=e([a,b])=f([a,b])=0\).
The boundary map
\[
 L_{r+1,s}\oplus L_{r,s+1}\longrightarrow L_{r+1,s+1},
 \qquad (u,v)\longmapsto[u,b]+[a,v],
\]
is surjective: a free Lie algebra generated by \(a,b\) is the
sum of \([a,L]\) and \([b,L]\) in every degree greater than one.
The multigraded Witt formula gives
\[
 \dim L_{3,2}=\dim L_{2,3}=2,\quad \dim L_{3,3}=3,
 \qquad
 \dim L_{4,3}=\dim L_{3,4}=5,\quad \dim L_{4,4}=8.
\]
Thus \(\dim D_{22}=1\) and \(\dim D_{33}=2\).
The displayed nonzero \(d\) spans \(D_{22}\);
the distinct last two \(b\)-coefficient vectors show that \(e,f\)
are independent and hence form a basis of \(D_{33}\).

To verify injectivity of the bracket map, only two associative
word coefficients are needed.  Expansion by the derivation rule
gives
\[
\begin{array}{c|rr}
 \text{coefficient in the image of }a &[d,e]&[d,f]\\ \hline
 aaaaabbabbb&-2&6\\
 aaaababbabb&-22&18
\end{array}
\qquad
\det\begin{pmatrix}-2&6\\-22&18\end{pmatrix}=96.
\]
Consequently \([d,\cdot]:D_{33}\longrightarrow D_{55}\) is
injective over \(\mathbb Q_p\), for every \(p\).
All equalities here are integral free-associative-algebra
identities.  The accompanying exact program
\texttt{check\_cm\_low\_class\_two.py} reconstructs the Lyndon
bases, checks the boundary equations, and records the complete
coefficient certificate.

We now identify the arithmetic classes to which this calculation
applies.  By \cite[Theorem 2.8 and Lemma 2.10, p.~9]{Ishii},
\(\kappa_{(2,2)}\) and \(\kappa_{(3,3)}\) are nonzero on the
weight-four and weight-six graded quotients, respectively.
Their nonvanishing uses the proven finiteness of the
second cohomology groups with Tate twists two and three.
The CM character decomposition of the rational graded pieces
is semisimple by \cite[Proposition 2.2]{Ishii}.
It therefore supplies nonzero rational classes of pure
bidegrees \((2,2)\) and \((3,3)\).
Clearing their denominators gives integral graded classes,
represented by actual elements \(h_4\in F^4H\) and
\(h_6\in F^6H\).
These elements are not asserted to satisfy exact eigen-equations
under conjugation.

The symplectic derivation comparison
\cite[Proposition 3.3(4), pp.~12--14]{IshiiLie} identifies these
two outer graded classes with nonzero elements of \(D_{22}\)
and \(D_{33}\).  Its inner kernel is generated by
\(\operatorname{ad}[a,b]\) in bidegree \((1,1)\);
it has zero component in bidegrees \((2,2),(3,3),(5,5)\).
The calculation therefore proves that
\([h_4,h_6]\) has nonzero outer class in weight ten.
In particular its image in \(\operatorname{Out}(P)\) is nontrivial.

By \cite[Proposition 2.2(1)]{Ishii}, the weight-one,
weight-two, and weight-three arithmetic graded pieces vanish.
Thus \(\rho_H(H)\subseteq F^4\operatorname{Out}(\Pi)\).
Centrality of the weight filtration gives
\[
 \rho_H(\gamma_3H)\subseteq F^{12}\operatorname{Out}(\Pi)
       =\ker\bigl(\operatorname{Out}(\Pi)\longrightarrow
                              \operatorname{Out}(\Pi/\Pi(13))\bigr).
\]
Every \(Q_r\) consequently has class at most two.

Finally, the natural map
\(\operatorname{Out}(P)\longrightarrow
 \varprojlim_r\operatorname{Out}(P_r)\) is injective.
To see this, suppose an automorphism is inner modulo every
power subgroup defining \(P_r\).  For each \(r\), its possible
conjugating elements form a nonempty closed subset of the
compact group \(P\), and these subsets are nested.
Their intersection contains a conjugator on \(P\), since
the intersection of those power subgroups is trivial.
The latter assertion follows because every finite quotient of
the pro-\(p\) group \(P\) has finite \(p\)-power exponent.
Hence the nontrivial outer image of \([h_4,h_6]\) survives
in some \(\operatorname{Out}(P_r)\) and in every later one.
This proves the proposition.
\end{proof}

The proposition gives an explicit nonzero central bracket in
the actual finite outer images.  Its source is the displayed
homomorphism \(H\to Q_r\); it does not identify the kernel of
the different homomorphism \(H\to H/\gamma_3H\).
